\section{Finite-Dimensional Vector Spaces}

\subsection{Span and Linear Independence}

\begin{exercise}[2a1]
    Find a list of four distinct vectors in $\f^3$ whose span equals
    \[\set{(x, y, z) \in \f^3 \mid x + y + z = 0}\]
\end{exercise}

\begin{sol}
    Clearly, the span of the list involves two vectors, since we may rewrite $z$ in terms of $x$ and $y$ as $z = -x - y$. Thus, we may choose the following two vectors:
    \[v_1 = (1, 0, -1), \quad v_2 = (0, 1, -1).\]
    To find two more vectors, we may pick any two distinct linear combinations of $v_1$ and $v_2$. For example, we may select
    \[v_3 = v_1 + v_2 = (1, 1, -2), \quad v_4 = 2v_1 + 3v_2 = (2, 3, -5).\]
    Then
    \[\spn(v_1, v_2, v_3, v_4) = \spn(v_1, v_2) = \set{(x, y, z) \in \f^3 \mid x + y + z = 0}. \qh\]
\end{sol}

\begin{exercise}[2a2]
    Prove or give a counterexample: If $v_1, v_2, v_3, v_4$ spans $V$, then the list
    \[v_1 - v_2, v_2 - v_3, v_3 - v_4, v_4\]
    also spans $V$.
\end{exercise}

\begin{sol}
    Let $U = \set{v_1, v_2, v_3, v_4}$ and $W = \set{v_1 - v_2, v_2 - v_3, v_3 - v_4, v_4}$. Clearly, $v_4 \in \spn U$. Then we may write
    \begin{align*}
        v_3 &= (v_3 - v_4) + v_4 \in \spn W, \\
        v_2 &= (v_2 - v_3) + v_3 \in \spn W, \\
        v_1 &= (v_1 - v_2) + v_2 \in \spn W.
    \end{align*}
    Thus, $v_1, v_2, v_3, v_4 \in \spn W$, and so $V = \spn U \subseteq \spn W$. Since $W$ is a list of vectors in $V$, we have $\spn W \subseteq V$. Therefore, $\spn W = V$, and so $W$ spans $V$.
\end{sol}
 
\begin{exercise}[2a3]
    Suppose $v_1, \ldots, v_m$ is a list of vectors in $V$. For $k \in \set{1, \ldots, m}$, let
    \[w_k = v_1 + \cdots + v_k.\]
    Show that $\spn(v_1, \ldots, v_m) = \spn(w_1, \ldots, w_m)$.
\end{exercise}

\begin{sol}
    This is a generalization of \autoref{ex:2a2} and a direct consequence of the fact that $v_k = w_k - w_{k-1}$ for $k \in \set{2, \ldots, m}$ and $v_1 = w_1$. Since each $v_k$ is a linear combination of $w_1, \ldots, w_k$, we have $\spn(v_1, \ldots, v_m) \subseteq \spn(w_1, \ldots, w_m)$. On the other hand, each $w_k$ is a linear combination of $v_1, \ldots, v_k$, so $\spn(w_1, \ldots, w_m) \subseteq \spn(v_1, \ldots, v_m)$. Therefore, $\spn(v_1, \ldots, v_m) = \spn(w_1, \ldots, w_m)$.
\end{sol}
 
\begin{exercise}[2a4]
    \begin{subexercises}
        \item Show that a list of length one in a vector space is linearly independent if and only if the vector in the list is not 0.
        \item Show that a list of length two in a vector space is linearly independent if and only if neither of the two vectors in the list is a scalar multiple of the other.
    \end{subexercises}
\end{exercise}

\begin{sole}
    \item Let $v$ be the vector in the list. If $v = 0$, then the list is linearly dependent, since $0 \cdot v = 0$. If $v \neq 0$, then the only way to write $0$ as a linear combination of $v$ is to take the scalar multiple to be $0$. Thus, the list is linearly independent.
    \item Let $v_1$ and $v_2$ be the vectors in the list. If $v_1$ is a scalar multiple of $v_2$, then there exists a scalar $\lambda$ such that $v_1 = \lambda v_2$. Then we may write
    \[0 = v_1 - \lambda v_2,\]
    which shows that the list is linearly dependent. Conversely, if the list is linearly dependent, then there exist scalars $\alpha$ and $\beta$, not both zero, such that
    \[\alpha v_1 + \beta v_2 = 0.\]
    If $\alpha \neq 0$, then
    \[v_1 = -\frac{\beta}{\alpha} v_2.\]
    Similarly, if $\beta \neq 0$, then
    \[v_2 = -\frac{\alpha}{\beta} v_1.\]
    Thus, the list is linearly independent if and only if neither vector is a scalar multiple of the other.
\end{sole}
 
\begin{exercise}[2a5]
    Find a number $t$ such that
    \[(3, 1, 4), (2, -3, 5), (5, 9, t)\]
    is not linearly independent in $\mathbb{R}^3$.
\end{exercise}

\begin{sol}
    To say that the list is not linearly independent is to say that there exist scalars $\alpha, \beta, \gamma$, not all zero, such that
    \[\alpha(3, 1, 4) + \beta(2, -3, 5) + \gamma(5, 9, t) = (0, 0, 0).\]
    This is equivalent to the following system of equations:
    \begin{align*}
        0 & = 3\alpha + 2\beta + 5\gamma, \\
        0 & = \alpha - 3\beta + 9\gamma, \\
        0 & = 4\alpha + 5\beta + t\gamma.
    \end{align*}
    From the second equation, we obtain $\alpha = 3\beta - 9\gamma$. Substituting this into the first equation, we have
    \begin{align*}
        0 & = 3(3\beta - 9\gamma) + 2\beta + 5\gamma \\
        & = 11\beta - 22\gamma.
    \end{align*}
    Thus, $\beta = 2\gamma$. Substituting this into the expression for $\alpha$, we have $\alpha = 3(2\gamma) - 9\gamma = -3\gamma$. Substituting these values into the third equation, we have
    \begin{align*}
        0 & = 4(-3\gamma) + 5(2\gamma) + t\gamma \\
        & = -12\gamma + 10\gamma + t\gamma \\
        & = (t - 2)\gamma.
    \end{align*}
    Since $\gamma \neq 0$, we must have $t - 2 = 0$, or $t = 2$. Then the list is not linearly independent when $t = 2$.
\end{sol}
 
\begin{exercise}[2a6]
    Show that the list $(2, 3, 1), (1, -1, 2), (7, 3, c)$ is linearly dependent in $\mathbb{F}^3$ if and only if $c = 8$.
\end{exercise}

\begin{sol}
    We may proceed similarly to \autoref{ex:2a5} to obtain the following system of equations:
    \begin{align*}
        0 & = 2\alpha + \beta + 7\gamma, \\
        0 & = 3\alpha - \beta + 3\gamma, \\
        0 & = \alpha + 2\beta + c\gamma.
    \end{align*}
    The second equation gives us $\beta = 3\alpha + 3\gamma$. Substituting this into the first equation, we have
    \begin{align*}
        0 & = 2\alpha + (3\alpha + 3\gamma) + 7\gamma \\
        & = 5\alpha + 10\gamma.
    \end{align*}
    Thus, $\alpha = -2\gamma$. Substituting this into the expression for $\beta$, we have $\beta = 3(-2\gamma) + 3\gamma = -3\gamma$. Substituting these values into the third equation, we have
    \begin{align*}
        0 & = (-2\gamma) + 2(-3\gamma) + c\gamma \\
        & = -2\gamma - 6\gamma + c\gamma \\
        & = (c - 8)\gamma.
    \end{align*}
    Since $\gamma \neq 0$, we must have $c - 8 = 0$, or $c = 8$. Then the list is linearly dependent when $c = 8$.

    Conversely, if $c \neq 8$, then $(c - 8)\gamma = 0$ implies $\gamma = 0$. Then $\alpha = -2\gamma = 0$ and $\beta = -3\gamma = 0$, so the only solution to the system of equations is $\alpha = \beta = \gamma = 0$. Thus, the list is linearly independent when $c \neq 8$.
\end{sol}
 
\begin{exercise}[2a7]
    \begin{subexercises}
        \item Show that if we think of $\c$ as a vector space over $\r$, then the list $1 + i, 1 - i$ is linearly independent.
        \item Show that if we think of $\c$ as a vector space over $\c$, then the list $1 + i, 1 - i$ is linearly dependent.
    \end{subexercises}
\end{exercise}

\begin{sole}
    \item Suppose $\alpha, \beta \in \r$ such that
    \[\alpha(1 + i) + \beta(1 - i) = 0.\]
    Then we have
    \[(\alpha + \beta) + (\alpha - \beta)i = 0.\]
    Since $0 = 0 + 0i$, we must have $\alpha + \beta = 0$ and $\alpha - \beta = 0$. Solving this system of equations, we find that $\alpha = \beta = 0$. Thus, the list is linearly independent.
    \item Suppose $a + bi, c + di \in \c$ such that
    \[(a + bi)(1 + i) + (c + di)(1 - i) = 0.\]
    Then we have
    \[(a - b + c + d) + (a + b - c + d)i = 0.\]
    Then we have the system of equations
    \begin{align*}
        0 & = a - b + c + d, \\
        0 & = a + b - c + d.
    \end{align*}
    Adding the two equations results in $0 = 2a + 2d$, or $a = -d$. Substituting this into the first equation, we obtain $b = c$. Then we have infinitely many nontrivial solutions to the system of equations, such as $a = 1, b = 0, c = 0, d = -1$. Thus, the list is linearly dependent.
\end{sole}

\begin{exercise}[2a8]
    Suppose $v_1, v_2, v_3, v_4$ is linearly independent in $V$. Prove that the list
    \(v_1 - v_2, v_2 - v_3, v_3 - v_4, v_4\)
    is also linearly independent.
\end{exercise}

\begin{sol}
    Let $\alpha, \beta, \gamma, \delta \in \f$ such that
    \[\alpha(v_1 - v_2) + \beta(v_2 - v_3) + \gamma(v_3 - v_4) + \delta v_4 = 0.\]
    Then we have
    \[\alpha v_1 + (-\alpha + \beta)v_2 + (-\beta + \gamma)v_3 + (-\gamma + \delta)v_4 = 0.\]
    Since $v_1, v_2, v_3, v_4$ is linearly independent, we must have
    \begin{align*}
        0 & = \alpha, \\
        0 & = -\alpha + \beta, \\
        0 & = -\beta + \gamma, \\
        0 & = -\gamma + \delta.
    \end{align*}
    From the first equation, we have $\alpha = 0$. Substituting this into the second equation gives us $\beta = 0$. Substituting this into the third equation gives us $\gamma = 0$. Finally, substituting this into the fourth equation gives us $\delta = 0$. Thus, the only solution to the original equation is $\alpha = \beta = \gamma = \delta = 0$, and so the list is linearly independent.
\end{sol}
 
\begin{exercise}[2a9]
    Prove or give a counterexample: If $v_1, v_2, \ldots, v_m$ is a linearly independent list of vectors in $V$, then
    \(5v_1 - 4v_2, v_2, v_3, \ldots, v_m\)
    is linearly independent.
\end{exercise}

\begin{sol}
    Let $\alpha_1, \ldots, \alpha_m \in \f$ such that
    \[\alpha_1(5v_1 - 4v_2) + \alpha_2 v_2 + \cdots + \alpha_m v_m = 0.\]
    Then we have
    \[5\alpha_1 v_1 + (-4\alpha_1 + \alpha_2)v_2 + \alpha_3 v_3 + \cdots + \alpha_m v_m = 0.\]
    Since $v_1, v_2, \ldots, v_m$ is linearly independent, we must have
    \begin{align*}
        0 & = 5\alpha_1, \\
        0 & = -4\alpha_1 + \alpha_2, \\
        0 & = \alpha_3, \\
        & \vdots \\
        0 & = \alpha_m.
    \end{align*}
    From the first equation, we have $\alpha_1 = 0$. Substituting this into the second equation gives us $\alpha_2 = 0$. The remaining equations give us $\alpha_k = 0$ for $k = 3, \ldots, m$. Thus, the only solution to the original equation is $\alpha_k = 0$ for $k = 1, \ldots, m$, and so the list is linearly independent.
\end{sol}

\begin{exercise}[2a10]
    Prove or give a counterexample: If $v_1, v_2, \ldots, v_m$ is a linearly independent list of vectors in $V$ and $\lambda \in \f$ with $\lambda \neq 0$, then $\lambda v_1, \lambda v_2, \ldots, \lambda v_m$ is linearly independent.
\end{exercise}

\begin{sol}
    Let $\alpha_1, \ldots, \alpha_m \in \f$ such that
    \[\alpha_1(\lambda v_1) + \cdots + \alpha_m(\lambda v_m) = 0.\]
    Then we have
    \[\lambda(\alpha_1 v_1 + \cdots + \alpha_m v_m) = 0.\]
    Since $\lambda \neq 0$, we must have
    \[\alpha_1 v_1 + \cdots + \alpha_m v_m = 0.\]
    Since $v_1, v_2, \ldots, v_m$ is linearly independent, we must have $\alpha_k = 0$ for $k = 1, \ldots, m$. Thus, the only solution to the original equation is $\alpha_k = 0$ for $k = 1, \ldots, m$, and so the list is linearly independent.
\end{sol}
 
\begin{exercise}[2a11]
    Prove or give a counterexample: If $v_1, \ldots, v_m$ and $w_1, \ldots, w_m$ are linearly independent lists of vectors in $V$, then the list $v_1 + w_1, \ldots, v_m + w_m$ is linearly independent.
\end{exercise}

\begin{sol}
    Consider the following counterexample in $\r^2$:
    \[v_1 = (1, 0), \quad v_2 = (0, 1), \quad w_1 = (0, 1), \quad w_2 = (1, 0).\]
    Then $v_1, v_2$ and $w_1, w_2$ are linearly independent lists of vectors in $\r^2$. However,
    \[v_1 + w_1 = (1, 1), \quad v_2 + w_2 = (1, 1),\]
    and so the list $v_1 + w_1, v_2 + w_2$ is linearly dependent. Thus, the statement is false.
\end{sol}
 
\begin{exercise}[2a12]
    Suppose $v_1, \ldots, v_m$ is linearly independent in $V$ and $w \in V$. Prove that if $v_1 + w, \ldots, v_m + w$ is linearly dependent, then $w \in \spn(v_1, \ldots, v_m)$.
\end{exercise}

\begin{sol}
    Let $\alpha_1, \ldots, \alpha_m \in \f$, not all zero, such that
    \[\alpha_1(v_1 + w) + \cdots + \alpha_m(v_m + w) = 0.\]
    Then we have
    \[(\alpha_1 v_1 + \cdots + \alpha_m v_m) + (\alpha_1 + \cdots + \alpha_m)w = 0.\]
    If $\alpha_1 + \cdots + \alpha_m = 0$, then we have
    \[\alpha_1 v_1 + \cdots + \alpha_m v_m = 0,\]
    which contradicts the linear independence of $v_1, \ldots, v_m$. Thus, we must have $\alpha_1 + \cdots + \alpha_m \neq 0$. Let $\beta = \alpha_1 + \cdots + \alpha_m$. Then we may write
    \[w = -\frac{\alpha_1}{\beta}v_1 - \cdots - \frac{\alpha_m}{\beta}v_m.\]
    Thus, $w$ is a linear combination of $v_1, \ldots, v_m$, and so $w \in \spn(v_1, \ldots, v_m)$.
\end{sol}
 
\begin{exercise}[2a13]
    Suppose $v_1, \ldots, v_m$ is linearly independent in $V$ and $w \in V$. Show that
    \[v_1, \ldots, v_m, w \text{ is linearly independent} \iff w \notin \spn(v_1, \ldots, v_m).\]
\end{exercise}

\begin{sol}
    Suppose $v_1, \ldots, v_m, w$ is linearly independent. If there were scalars $\alpha_1, \ldots, \alpha_m, \beta$ such that
    \[\alpha_1 v_1 + \cdots + \alpha_m v_m + \beta w = 0,\]
    then linear independence would imply that $\alpha_1 = \cdots = \alpha_m = \beta = 0$. Suppose $w \in \spn(v_1, \ldots, v_m)$. Then there exist scalars $\gamma_1, \ldots, \gamma_m$ not all zero such that
    \[\gamma_1 v_1 + \cdots + \gamma_m v_m = w.\]
    Then we may write
    \[-\gamma_1 v_1 - \cdots - \gamma_m v_m + 1\cdot w = 0,\]
    which contradicts the linear independence of $v_1, \ldots, v_m, w$. Thus, $w \notin \spn(v_1, \ldots, v_m)$.

    Conversely, suppose $w \notin \spn(v_1, \ldots, v_m)$. Then there do not exist scalars $\alpha_1, \ldots, \alpha_m$ such that
    \[\alpha_1 v_1 + \cdots + \alpha_m v_m = w.\]
    If $v_1, \ldots, v_m, w$ is linearly dependent, then there exist scalars $\beta_1, \ldots, \beta_m, \gamma$, not all zero, such that
    \[\beta_1 v_1 + \cdots + \beta_m v_m + \gamma w = 0.\]
    If $\gamma = 0$, then we have
    \[\beta_1 v_1 + \cdots + \beta_m v_m = 0,\]
    which contradicts the linear independence of $v_1, \ldots, v_m$. Thus, we must have $\gamma \neq 0$. Then we may write
    \[w = -\frac{\beta_1}{\gamma}v_1 - \cdots -\frac{\beta_m}{\gamma}v_m,\]
    which contradicts the assumption that $w$ is not in the span of $v_1, \ldots, v_m$. Thus, $v_1, \ldots, v_m, w$ is linearly independent.
\end{sol}
 
\begin{exercise}[2a14]
    Suppose $v_1, \ldots, v_m$ is a list of vectors in $V$. For $k \in \set{1, \ldots, m}$, let
    \[w_k = v_1 + \cdots + v_k.\]
    Show that the list $v_1, \ldots, v_m$ is linearly independent if and only if the list $w_1, \ldots, w_m$ is linearly independent.
\end{exercise}

\begin{sol}
    Suppose $v_1, \ldots, v_m$ is linearly independent. Let $\alpha_1, \ldots, \alpha_m \in \f$ such that
    \[\alpha_1 w_1 + \cdots + \alpha_m w_m = 0.\]
    Noting that $v_k = w_k - w_{k-1}$ for $k \in \set{2, \ldots, m}$ and $v_1 = w_1$, we may rewrite the equation as
    \[\alpha_1 v_1 + (\alpha_2 - \alpha_1)v_2 + \cdots + (\alpha_m - \alpha_{m-1})v_m = 0.\]
    Since $v_1, \ldots, v_m$ is linearly independent, we must have $\alpha_1 = 0$. Then this implies that $\alpha_2 - \alpha_1 = 0$, and so on, until we have $\alpha_m - \alpha_{m-1} = 0$. Thus, the only solution to the original equation is $\alpha_1 = \cdots = \alpha_m = 0$, and so $w_1, \ldots, w_m$ is linearly independent.

    Conversely, suppose $w_1, \ldots, w_m$ is linearly independent. Let $\beta_1, \ldots, \beta_m \in \f$ such that
    \[\beta_1 v_1 + \cdots + \beta_m v_m = 0.\]
    Then we have
    \[\beta_1 w_1 + (\beta_2 + \beta_1) w_2 + \cdots + (\beta_m + \cdots + \beta_1) w_m = 0.\]
    Since $w_1, \ldots, w_m$ is linearly independent, we must have $\beta_1 = 0$. Then this implies that $\beta_2 + \beta_1 = 0$, and so on, until we have $\beta_m + \cdots + \beta_1 = 0$. Thus, the only solution to the original equation is $\beta_1 = \cdots = \beta_m = 0$, and so $v_1, \ldots, v_m$ is linearly independent.
\end{sol}
 
\begin{exercise}[2a15]
    Explain why there does not exist a list of six polynomials that is linearly independent in $\pp_4(\f)$.
\end{exercise}

\begin{sol}
    Since $\pp_4(\f)$ can be spanned by the list $1, x, x^2, x^3, x^4$, and the length of a linearly independent list cannot exceed the length of a spanning list, there does not exist a list of six polynomials that is linearly independent in $\pp_4(\f)$.
\end{sol}
 
\begin{exercise}[2a16]
    Explain why no list of four polynomials spans $\pp_4(\f)$.
\end{exercise}

\begin{sol}
    Since $\pp_4(\f)$ can be spanned by the list $1, x, x^2, x^3, x^4$, and the length of a spanning list cannot be less than the length of a linearly independent list, no list of four polynomials spans $\pp_4(\f)$.
\end{sol}
 
\begin{exercise}[2a17]
    Prove that $V$ is infinite-dimensional if and only if there is a sequence $v_1, v_2, \ldots$ of vectors in $V$ such that $v_1, \ldots, v_m$ is linearly independent for every positive integer $m$.
\end{exercise}

\begin{sol}
    Suppose $V$ is infinite-dimensional. We proceed by induction to construct a sequence of vectors $v_1, v_2, \ldots$ in $V$ such that $v_1, \ldots, v_m$ is linearly independent for every positive integer $m$. Since $V$ is infinite-dimensional, there exists a nonzero vector $v_1 \in V$. Suppose we have constructed a linearly independent list of vectors $v_1, \ldots, v_m$ in $V$. Since $V$ is infinite-dimensional, there exists a vector $v_{m+1} \in V$ such that $v_{m + 1} \notin \spn(v_1, \ldots, v_m)$. Then $v_1, \ldots, v_m, v_{m+1}$ is linearly independent. By induction, we have constructed a sequence of vectors $v_1, v_2, \ldots$ in $V$ such that $v_1, \ldots, v_m$ is linearly independent for every positive integer $m$.

    Conversely, suppose there exists a sequence of vectors $v_1, v_2, \ldots$ in $V$ such that $v_1, \ldots, v_m$ is linearly independent for every positive integer $m$. Then for every positive integer $m$, there exists a linearly independent list of vectors in $V$ of length $m$. Since there is no maximum length for a linearly independent list of vectors in $V$, this implies that $V$ is infinite-dimensional.
\end{sol}
 
\begin{exercise}[2a18]
    Prove that $\f^\infty$ is infinite-dimensional.
\end{exercise}

\begin{sol}
    By \autoref{ex:2a17}, it suffices to construct a sequence of vectors $v_1, v_2, \ldots$ in $\f^\infty$ such that $v_1, \ldots, v_m$ is linearly independent for every positive integer $m$. Let $v_k$ be the vector in $\f^\infty$ whose $k$-th coordinate is $1$ and all other coordinates are $0$. Explicitly, we have
    \[v_1 = (1, 0, 0, \ldots), \quad v_2 = (0, 1, 0, \ldots), \quad v_3 = (0, 0, 1, \ldots), \quad \ldots\]
    Then for any positive integer $m$, the list $v_1, \ldots, v_m$ is linearly independent, since the only way to write the zero vector as a linear combination of $v_1, \ldots, v_m$ is to take all coefficients to be zero. Thus, by \autoref{ex:2a17}, $\f^\infty$ is infinite-dimensional.
\end{sol}
 
\begin{exercise}[2a19]
    Prove that the real vector space of all continuous real-valued functions on the interval $[0, 1]$ is infinite-dimensional.
\end{exercise}

\begin{sol}
    By \autoref{ex:2a17}, it suffices to construct a sequence of continuous real-valued functions $f_1, f_2, \ldots$ on the interval $[0, 1]$ such that $f_1, \ldots, f_m$ is linearly independent for every positive integer $m$. Let $f_k(x) = x^{k-1}$ for $k = 1, 2, \ldots$. Then for any positive integer $m$, the list $f_1, \ldots, f_m$ is linearly independent, since the only way to write the zero function as a linear combination of $f_1, \ldots, f_m$ is to take all coefficients to be zero. Thus, by \autoref{ex:2a17}, $\r^{[0, 1]}$ is infinite-dimensional.
\end{sol}
 
\begin{exercise}[2a20]
    Suppose $p_0, p_1, \ldots, p_m$ are polynomials in $\pp_m(\f)$ such that $p_k(2) = 0$ for each $k \in \set{0, \ldots, m}$. Prove that $p_0, p_1, \ldots, p_m$ is not linearly independent in $\pp_m(\f)$.
\end{exercise}

\begin{sol}
    Since $p_k(2) = 0$ for each $k \in \set{0, \ldots, m}$, we have
    \[p_k(x) = (x - 2)q_k(x)\]
    for some polynomial $q_k(x) \in \pp_{m-1}(\f)$. Then we may write
    \[p_0(x), p_1(x), \ldots, p_m(x) = (x - 2)q_0(x), (x - 2)q_1(x), \ldots, (x - 2)q_m(x).\]
    Since $\pp_{m - 1}(\f)$ can be spanned by the $m$ polynomials $1, x, \ldots, x^{m-1}$, the list of $m + 1$ polynomials $q_0, q_1, \ldots, q_m$ must be linearly dependent. Thus, there exist scalars $\alpha_0, \ldots, \alpha_m$, not all zero, such that
    \[\alpha_0 q_0(x) + \cdots + \alpha_m q_m(x) = 0.\]
    Then we may write
    \[\alpha_0 p_0(x) + \cdots + \alpha_m p_m(x) = (x - 2)(\alpha_0 q_0(x) + \cdots + \alpha_m q_m(x)) = 0.\]
    Since not all the scalars $\alpha_0, \ldots, \alpha_m$ are zero, this shows that $p_0, p_1, \ldots, p_m$ is linearly dependent in $\pp_m(\f)$.
\end{sol}

\newpage

\subsection{Bases}

\begin{exercise}[2b1]
    Find all vector spaces that have exactly one basis.
\end{exercise}

\begin{sol}
    Let $V$ be a vector space. If $V$ has exactly one basis, then $V$ must be finite-dimensional, since infinite-dimensional vector spaces have infinitely many bases. Let $v_1, \ldots, v_n$ be a basis of $V$. Then any other basis of $V$ must also have length $n$. Since there is only one basis of $V$, it must be that $n = 0$, and so the only vector space that has exactly one basis is the zero vector space.
\end{sol}

\begin{exercise}[2b2]
    Verify all assertions in Example 2.27.
    \begin{subexercises}
        \item The list $(1, 0, \ldots, 0), (0, 1, 0, \ldots, 0), \ldots, (0, \ldots, 0, 1)$ is a basis of $\f^n$, called the \emph{standard basis} of $\f^n$.
        \item The list $(1, 2), (3, 5)$ is a basis of $\f^2$. Note that this list has length two, which is the same as the length of the standard basis of $\f^2$. In the next section, we will see that this is not a coincidence.
        \item The list $(1, 2, -4), (7, -5, 6)$ is linearly independent in $\f^3$ but is not a basis of $\f^3$ because it does not span $\f^3$.
        \item The list $(1, 2), (3, 5), (4, 13)$ spans $\f^2$ but is not a basis of $\f^2$ because it is not linearly independent.
        \item The list $(1, 1, 0), (0, 0, 1)$ is a basis of $\set{(x, x, y) \in \f^3 \mid x, y \in \f}$.
        \item The list $(1, -1, 0), (1, 0, -1)$ is a basis of
        \[\set{(x, y, z) \in \f^3 \mid x + y + z = 0}.\]
        \item The list $1, z, \ldots, z^m$ is a basis of $\pp_m(\f)$, called the \emph{standard basis} of $\pp_m(\f)$.
    \end{subexercises}
\end{exercise}

\begin{sole}
    \item Suppose $\alpha_1, \ldots, \alpha_n \in \f$ such that
    \[\alpha_1(1, 0, \ldots, 0) + \alpha_2(0, 1, 0, \ldots, 0) + \cdots + \alpha_n(0, \ldots, 0, 1) = 0.\]
    Then each $\alpha_i = 0$, so the list is linearly independent. Moreover, any vector $(x_1, \ldots, x_n) \in \f^n$ can be written as
    \[(x_1, \ldots, x_n) = x_1(1, 0, \ldots, 0) + x_2(0, 1, 0, \ldots, 0) + \cdots + x_n(0, \ldots, 0, 1),\]
    so the list spans $\f^n$. Thus, the list is a basis of $\f^n$.
    \item If $(3, 5) = a(1, 2)$ for some $a \in \f$, then we must have $3 = a$ and $5 = 2a$, which is impossible. Thus, the list is linearly independent. Moreover, for any element $(x, y) \in \f^2$, we can solve the system of equations
    \begin{align*}
        x & = \alpha + 3\beta, \\
        y & = 2\alpha + 5\beta
    \end{align*}
    for $\alpha, \beta \in \f$. Solving this system of equations, we find that
    \[\alpha = 5x - 3y, \quad \beta = -2x + y.\]
    Since we can find $\alpha, \beta \in \f$ for any $(x, y) \in \f^2$, the list spans $\f^2$. Thus, the list is a basis of $\f^2$.
    \item If $(7, -5, 6) = a(1, 2, -4)$ for some $a \in \f$, then we must have $7 = a$, $-5 = 2a$, and $6 = -4a$, which is impossible. Thus, the list is linearly independent. Moreover, the list does not span $\f^3$. If it did, then we could write $(0, 0, 1) = \alpha(1, 2, -4) + \beta(7, -5, 6)$ for some $\alpha, \beta \in \f$. This would give us the system of equations
    \begin{align*}
        0 & = \alpha + 7\beta, \\
        0 & = 2\alpha - 5\beta, \\
        1 & = -4\alpha + 6\beta.
    \end{align*}
    From the first equation, we see that $\alpha = -7\beta$. Substituting this into the second equation gives us $0 = -14\beta - 5\beta = -19\beta$, so $\beta = 0$ and $\alpha = 0$. However, substituting these values into the third equation gives us $1 = 0$, which is impossible. Thus, the list does not span $\f^3$, and so it is not a basis of $\f^3$.
    \item From part (b), we saw that $(1, 2), (3, 5)$ is a basis of $\f^2$. Since $(4, 13) = 5(1, 2) - 2(3, 5)$, the list $(1, 2), (3, 5), (4, 13)$ is linearly dependent. Moreover, since $(1, 2), (3, 5)$ is a basis of $\f^2$, the list $(1, 2), (3, 5), (4, 13)$ spans $\f^2$. Thus, the list is not a basis of $\f^2$.
    \item Let $U$ be the subspace in question. If $\alpha, \beta \in \f$ such that
    \[\alpha(1, 1, 0) + \beta(0, 0, 1) = (0, 0, 0),\]
    then we must have $\alpha = \beta = 0$, so the list is linearly independent. Moreover, for any $(x, x, y) \in U$, we can write
    \[(x, x, y) = x(1, 1, 0) + y(0, 0, 1),\]
    so the list spans $U$. Thus, the list is a basis of $U$.
    \item Let $U$ be the subspace in question. If $\alpha, \beta \in \f$ such that
    \[\alpha(1, -1, 0) + \beta(1, 0, -1) = (0, 0, 0),\]
    then we must have $\alpha = \beta = 0$, so the list is linearly independent. Moreover, for any $(x, y, z) \in U$, we can write
    \[(x, y, z) = (x - z)(1, -1, 0) + (y - z)(1, 0, -1),\]
    so the list spans $U$. Thus, the list is a basis of $U$.
    \item If $\alpha_0, \ldots, \alpha_m \in \f$ such that
    \[\alpha_0 + \alpha_1 z + \cdots + \alpha_m z^m = 0,\]
    then we must have $\alpha_0 = \cdots = \alpha_m = 0$, so the list is linearly independent. Moreover, any polynomial $p(z) \in \pp_m(\f)$ can be written as
    \[p(z) = \alpha_0 + \alpha_1 z + \cdots + \alpha_m z^m,\]
    so the list spans $\pp_m(\f)$. Thus, the list is a basis of $\pp_m(\f)$.
\end{sole}

\begin{exercise}[2b3]
    \begin{subexercises}
        \item Let $U$ be the subspace of $\r^5$ defined by
        \[U = \set{(x_1, x_2, x_3, x_4, x_5) \in \r^5 \mid x_1 = 3x_2 \text{ and } x_3 = 7x_4}.\]
        Find a basis of $U$.
        \item Extend the basis in (a) to a basis of $\r^5$.
        \item Find a subspace $W$ of $\r^5$ such that $\r^5 = U \op W$.
    \end{subexercises}
\end{exercise}

\begin{sole}
    \item Note that we can rewrite the subspace $U$ as
    \[U = \set{(3x_2, x_2, 7x_4, x_4, x_5) \in \r^5 \mid x_2, x_4, x_5 \in \r}.\]
    Then we can write any vector in $U$ as a linear combination of the following three vectors:
    \[u_1 = (3, 1, 0, 0, 0), \quad u_2 = (0, 0, 7, 1, 0), \quad u_3 = (0, 0, 0, 0, 1).\]
    Moreover, these three vectors are linearly independent, since the only way to write the zero vector as a linear combination of $u_1, u_2, u_3$ is to take all coefficients to be zero. Thus, $\set{u_1, u_2, u_3}$ is a basis of $U$.
    \item We can extend the basis $\set{u_1, u_2, u_3}$ of $U$ to a basis of $\r^5$ by adding two more vectors that are linearly independent of $u_1, u_2, u_3$. For example, we can add the vectors
    \[u_4 = (1, 0, 0, 0, 0), \quad u_5 = (0, 0, 1, 0, 0),\]
    which are linearly independent of $u_1, u_2, u_3$ since they are not in $U$. Thus, $\set{u_1, u_2, u_3, u_4, u_5}$ is a basis of $\r^5$.
    \item Let $W$ be the subspace of $\r^5$ spanned by the vectors $u_4$ and $u_5$. Then we have
    \[\r^5 = U \op W,\]
    since $U \cap W = \set{0}$ and every vector in $\r^5$ can be written as a sum of a vector in $U$ and a vector in $W$.
\end{sole}

\begin{exercise}[2b4]
    \begin{subexercises}
        \item Let $U$ be the subspace of $\c^5$ defined by
        \[U = \set{(z_1, z_2, z_3, z_4, z_5) \in \c^5 \mid 6z_1 = z_2 \text{ and } z_3 + 2z_4 + 3z_5 = 0}.\]
        Find a basis of $U$.
        \item Extend the basis in (a) to a basis of $\c^5$.
        \item Find a subspace $W$ of $\c^5$ such that $\c^5 = U \op W$.
    \end{subexercises}
\end{exercise}

\begin{sole}
    \item Note that we can rewrite the subspace $U$ as
    \[U = \set{(z_1, 6z_1, -2z_4 - 3z_5, z_4, z_5) \in \c^5 \mid z_1, z_4, z_5 \in \c}.\]
    Then we can write any vector in $U$ as a linear combination of the following three vectors:
    \[u_1 = (1, 6, 0, 0, 0), \quad u_2 = (0, 0, -2, 1, 0), \quad u_3 = (0, 0, -3, 0, 1).\]
    Moreover, these three vectors are linearly independent, since the only way to write the zero vector as a linear combination of $u_1, u_2, u_3$ is to take all coefficients to be zero. Thus, $\set{u_1, u_2, u_3}$ is a basis of $U$.
    \item We can extend the basis $\set{u_1, u_2, u_3}$ of $U$ to a basis of $\c^5$ by adding two more vectors that are linearly independent of $u_1, u_2, u_3$. For example, we can add the vectors
    \[u_4 = (1, 0, 0, 0, 0), \quad u_5 = (0, 0, 1, 0, 0),\]
    which are linearly independent of $u_1, u_2, u_3$ since they are not in $U$. Thus, $\set{u_1, u_2, u_3, u_4, u_5}$ is a basis of $\c^5$.
    \item Let $W$ be the subspace of $\c^5$ spanned by the vectors $u_4$ and $u_5$. Then we have
    \[\c^5 = U \op W,\]
    since $U \cap W = \set{0}$ and every vector in $\c^5$ can be written as a sum of a vector in $U$ and a vector in $W$.
\end{sole}

\begin{exercise}[2b5]
    Suppose $V$ is finite-dimensional and $U, W$ are subspaces of $V$ such that $V = U + W$. Prove that there exists a basis of $V$ consisting of vectors in $U \cup W$.
\end{exercise}

\begin{sol}
    Let $u_1, \ldots, u_m$ be a basis of $U$ and $w_1, \ldots, w_n$ be a basis of $W$. Then the list $u_1, \ldots, u_m, w_1, \ldots, w_n$ spans $V$, since any vector in $V$ can be written as a sum of a vector in $U$ and a vector in $W$. We can extract a basis of $V$ from this list by removing any linearly dependent vectors. Thus, there exists a basis of $V$ consisting of vectors in $U \cup W$.
\end{sol}

\begin{exercise}[2b6]
    Prove or give a counterexample: If $p_0, p_1, p_2, p_3$ is a list in $\pp_3(\f)$ such that none of the polynomials $p_0, p_1, p_2, p_3$ has degree 2, then the list is not a basis of $\pp_3(\f)$.
\end{exercise}

\begin{sol}
    This is not true. Let $p_0 = 1, p_1 = x, p_2 = x^3, p_3 = x^2 + x^3$. Then none of the polynomials $p_0, p_1, p_2, p_3$ has degree 2. Moreover, the list $p_0, p_1, p_2, p_3$ is linearly independent, since the only way to write the zero polynomial as a linear combination of $p_0, p_1, p_2, p_3$ is to take all coefficients to be zero. Finally, noting that $x^2 = p_3 - p_2$, then for any polynomial $p(x) = a_0 + a_1 x + a_2 x^2 + a_3 x^3 \in \pp_3(\f)$, we can write
    \[p(x) = a_0 p_0 + a_1 p_1 + a_2 (p_3 - p_2) + a_3 p_2 = a_0 p_0 + a_1 p_1 + (a_3 - a_2)p_2 + a_2 p_3,\]
    which shows that the list spans $\pp_3(\f)$. Thus, the list $p_0, p_1, p_2, p_3$ is a basis of $\pp_3(\f)$.
\end{sol}

\begin{exercise}[2b7]
    Suppose $v_1, v_2, v_3, v_4$ is a basis of $V$. Prove that
    \[v_1 + v_2, v_2 + v_3, v_3 + v_4, v_4\]
    is also a basis of $V$.
\end{exercise}

\begin{sol}
    Let $\alpha_1, \alpha_2, \alpha_3, \alpha_4 \in \f$ such that
    \[\alpha_1(v_1 + v_2) + \alpha_2(v_2 + v_3) + \alpha_3(v_3 + v_4) + \alpha_4 v_4 = 0.\]
    Then we may rewrite this equation as
    \[\alpha_1 v_1 + (\alpha_1 + \alpha_2)v_2 + (\alpha_2 + \alpha_3)v_3 + (\alpha_3 + \alpha_4)v_4 = 0.\]
    Since $v_1, v_2, v_3, v_4$ is a basis of $V$, it is linearly independent, and so we must have
    \[\alpha_1 = 0, \quad \alpha_1 + \alpha_2 = 0, \quad \alpha_2 + \alpha_3 = 0, \quad \alpha_3 + \alpha_4 = 0.\]
    From the first equation, we have $\alpha_1 = 0$. Then the second equation gives us $\alpha_2 = 0$. The third equation then gives us $\alpha_3 = 0$, and finally the fourth equation gives us $\alpha_4 = 0$. Thus, the list is linearly independent. Moreover, since $v_1, v_2, v_3, v_4$ is a basis of $V$, it spans $V$, and so the list also spans $V$. Therefore, the list is a basis of $V$.
\end{sol}

\begin{exercise}[2b8]
    Prove or give a counterexample: If $v_1, v_2, v_3, v_4$ is a basis of $V$ and $U$ is a subspace of $V$ such that $v_1, v_2 \in U$ and $v_3 \notin U$ and $v_4 \notin U$, then $v_1, v_2$ is a basis of $U$.
\end{exercise}

\begin{sol}
    This is not true. Let $V = \r^4$ and $v_1, v_2, v_3, v_4$ be the standard basis of $\r^4$. Let $U$ be the subspace of $\r^4$ defined by
    \[U = \spn(v_1, v_2, v_3 + v_4) = \set{(x, y, z, z) \in \r^4 \mid x, y, z \in \r}.\]
    Then $v_1, v_2 \in U$ and $v_3, v_4 \notin U$. However, $v_1, v_2$ does not span $U$, since $(0, 0, 1, 1) \in U$ but cannot be written as a linear combination of $v_1$ and $v_2$. Thus, $v_1, v_2$ is not a basis of $U$.
\end{sol}

\begin{exercise}[2b9]
    Suppose $v_1, \ldots, v_m$ is a list of vectors in $V$. For $k \in \set{1, \ldots, m}$, let
    \[w_k = v_1 + \cdots + v_k.\]
    Show that $v_1, \ldots, v_m$ is a basis of $V$ if and only if $w_1, \ldots, w_m$ is a basis of $V$.
\end{exercise}

\begin{sol}
    By \autoref{ex:2a14}, we know that $v_1, \ldots, v_m$ is linearly independent if and only if $w_1, \ldots, w_m$ is linearly independent. By \autoref{ex:2a3}, we know that $v_1, \ldots, v_m$ spans $V$ if and only if $w_1, \ldots, w_m$ spans $V$. Thus, $v_1, \ldots, v_m$ is a basis of $V$ if and only if $w_1, \ldots, w_m$ is a basis of $V$.
\end{sol}

\begin{exercise}[2b10]
    Suppose $U$ and $W$ are subspaces of $V$ such that $V = U \op W$. Suppose also that $u_1, \ldots, u_m$ is a basis of $U$ and $w_1, \ldots, w_n$ is a basis of $W$. Prove that
    \[u_1, \ldots, u_m, w_1, \ldots, w_n\]
    is a basis of $V$.
\end{exercise}

\begin{sol}
    Let $u_1, \ldots, u_m$ be a basis of $U$ and $w_1, \ldots, w_n$ be a basis of $W$. We claim that the list
    \[u_1, \ldots, u_m, w_1, \ldots, w_n\]
    is a basis of $V$. Let $\alpha_1, \ldots, \alpha_m, \beta_1, \ldots, \beta_n \in \f$ such that
    \[\alpha_1 u_1 + \cdots + \alpha_m u_m + \beta_1 w_1 + \cdots + \beta_n w_n = 0.\]
    Then we may write
    \[\alpha_1 u_1 + \cdots + \alpha_m u_m = -\beta_1 w_1 - \cdots - \beta_n w_n.\]
    Since the left-hand side is in $U$ and the right-hand side is in $W$, and $U \cap W = \set{0}$, we must have
    \[\alpha_1 u_1 + \cdots + \alpha_m u_m = 0 \quad \text{and} \quad \beta_1 w_1 + \cdots + \beta_n w_n = 0.\]
    Since $u_1, \ldots, u_m$ is a basis of $U$, we have $\alpha_1 = \cdots = \alpha_m = 0$. Similarly, since $w_1, \ldots, w_n$ is a basis of $W$, we have $\beta_1 = \cdots = \beta_n = 0$. Thus, the list is linearly independent. Moreover, since $V = U \op W$, then $v = u + w$ for some $u \in U$ and $w \in W$. Then we may write
    \[u = \gamma_1 u_1 + \cdots + \gamma_m u_m \quad \text{and} \quad w = \delta_1 w_1 + \cdots + \delta_n w_n\]
    for some $\gamma_1, \ldots, \gamma_m, \delta_1, \ldots, \delta_n \in \f$. Hence,
    \[v = u + w = \gamma_1 u_1 + \cdots + \gamma_m u_m + \delta_1 w_1 + \cdots + \delta_n w_n,\]
    which shows that the list spans $V$. Therefore, the list is a basis of $V$.
\end{sol}

\begin{exercise}[2b11]
    Suppose $V$ is a real vector space. Show that if $v_1, \ldots, v_n$ is a basis of $V$ (as a real vector space), then $v_1, \ldots, v_n$ is also a basis of the complexification $V_\c$ (as a complex vector space).

    \note{Please see \autoref{ex:1b8} for the definition of the complexification $V_\c$}.
\end{exercise}

\begin{sol}
    Let $v_1, \ldots, v_n$ be a basis of $V$. We claim that $v_1, \ldots, v_n$ is also a basis of $V_\c$. Let $\alpha_1, \ldots, \alpha_n \in \c$ such that
    \[\alpha_1 v_1 + \cdots + \alpha_n v_n = 0.\]
    Then we may write each $\alpha_k = a_k + b_k i$ for some $a_k, b_k \in \r$. Then we have
    \[(a_1 + b_1 i)v_1 + \cdots + (a_n + b_n i)v_n = 0.\]
    This implies that
    \[(a_1 v_1 + \cdots + a_n v_n) + i(b_1 v_1 + \cdots + b_n v_n) = 0.\]
    Since the left-hand side is a sum of two vectors in $V$, and $V$ is a real vector space, we must have
    \[a_1 v_1 + \cdots + a_n v_n = 0 \quad \text{and} \quad b_1 v_1 + \cdots + b_n v_n = 0.\]
    Since $v_1, \ldots, v_n$ is a basis of $V$, we have $a_1 = \cdots = a_n = 0$ and $b_1 = \cdots = b_n = 0$. Thus, $\alpha_1 = \cdots = \alpha_n = 0$, which shows that the list is linearly independent. Moreover, let $v \in V_\c$. Then we may write $v = u + w i$ for some $u, w \in V$. Since $v_1, \ldots, v_n$ is a basis of $V$, we may write
    \[u = \gamma_1 v_1 + \cdots + \gamma_n v_n \quad \text{and} \quad w = \delta_1 v_1 + \cdots + \delta_n v_n\]
    for some $\gamma_1, \ldots, \gamma_n, \delta_1, \ldots, \delta_n \in \r$. Hence,
    \[v = u + w i = (\gamma_1 + \delta_1 i)v_1 + \cdots + (\gamma_n + \delta_n i)v_n,\]
    which shows that the list spans $V_\c$. Therefore, the list is a basis of $V_\c$.
\end{sol}

\newpage

\subsection{Dimension}

\begin{exercise}[2c1]
    Show that the subspaces of $\r^2$ are precisely $\set{0}$, all lines in $\r^2$ containing the origin, and $\r^2$.
\end{exercise}

\begin{sol}
    Since $\dim(\r^2) = 2$, any subspace of $\r^2$ must have dimension $0$, $1$, or $2$. The only subspace of dimension $0$ is $\set{0}$, and the only subspace of dimension $2$ is $\r^2$. Any subspace with dimension 1 consists of all scalar multiples of some nonzero $(x, y) \in \r^2$, which is precisely the line with direction $(x, y)$ containing the origin. Thus, the subspaces of $\r^2$ are precisely $\set{0}$, all lines in $\r^2$ containing the origin, and $\r^2$.
\end{sol}

\begin{exercise}[2c2]
    Show that the subspaces of $\r^3$ are precisely $\set{0}$, all lines in $\r^3$ containing the origin, all planes in $\r^3$ containing the origin, and $\r^3$.
\end{exercise}

\begin{sol}
    Since $\dim(\r^3) = 3$, any subspace of $\r^3$ must have dimension $0$, $1$, $2$, or $3$. The only subspace of dimension $0$ is $\set{0}$, and the only subspace of dimension $3$ is $\r^3$. Any subspace with dimension 1 consists of all scalar multiples of some nonzero $(x, y, z) \in \r^3$, which is precisely the line containing the origin in the direction of $(x, y, z)$. Any subspace with dimension 2 consists of all linear combinations of two linearly independent vectors in $\r^3$, which is precisely a plane containing the origin. Thus, the subspaces of $\r^3$ are precisely $\set{0}$, all lines in $\r^3$ containing the origin, all planes in $\r^3$ containing the origin, and $\r^3$.
\end{sol}

\begin{exercise}[2c3]
    \begin{subexercises}
        \item Let $U = \set{p \in \pp_4(\f) \mid p(6) = 0}$. Find a basis of $U$.
        \item Extend the basis in (a) to a basis of $\pp_4(\f)$.
        \item Find a subspace $W$ of $\pp_4(\f)$ such that $\pp_4(\f) = U \op W$.
    \end{subexercises}
\end{exercise}

\begin{sole}
    \item Let $p = a_0 + a_1 x + a_2 x^2 + a_3 x^3 + a_4 x^4 \in U$. Then we have $p(6) = 0$, which gives us the equation
    \[a_0 + 6a_1 + 36a_2 + 216a_3 + 1296a_4 = 0.\]
    We can solve this equation for $a_0$ in terms of $a_1, a_2, a_3, a_4$:
    \[a_0 = -6a_1 - 36a_2 - 216a_3 - 1296a_4.\]
    Then we can write any polynomial in $U$ as a linear combination of the following four polynomials:
    \[u_1 = -6 + x, \quad u_2 = -36 + x^2, \quad u_3 = -216 + x^3, \quad u_4 = -1296 + x^4.\]
    Moreover, these four polynomials are linearly independent, since the only way to write the zero polynomial as a linear combination of $u_1, u_2, u_3, u_4$ is to take all coefficients to be zero. Thus, $\set{u_1, u_2, u_3, u_4}$ is a basis of $U$.
    \item We can extend the basis $\set{u_1, u_2, u_3, u_4}$ of $U$ to a basis of $\pp_4(\f)$ by adding one more polynomial that is linearly independent of $u_1, u_2, u_3, u_4$. For example, we can add the polynomial
    \[u_5 = 1,\]
    which is linearly independent of $u_1, u_2, u_3, u_4$ since it is not in $U$. Thus, $\set{u_1, u_2, u_3, u_4, u_5}$ is a basis of $\pp_4(\f)$.
    \item Let $W = \spn u_5 = \spn 1$. Then we have
    \[\pp_4(\f) = U + W,\]
    since every polynomial in $\pp_4(\f)$ can be uniquely written as a sum of a polynomial in $U$ and a polynomial in $W$. Moreover, if $p \in U \cap W$, then $p = \alpha$ for some $\alpha \in \f$. Since $p(6) = 0$, we have $\alpha = 0$. Thus, $U \cap W = \set{0}$, and so we have
    \[\pp_4(\f) = U \op W. \qh\]
\end{sole}

\begin{exercise}[2c4]
    \begin{subexercises}
        \item Let $U = \set{p \in \pp_4(\r) \mid p''(6) = 0}$. Find a basis of $U$.
        \item Extend the basis in (a) to a basis of $\pp_4(\r)$.
        \item Find a subspace $W$ of $\pp_4(\r)$ such that $\pp_4(\r) = U \op W$.
    \end{subexercises}
\end{exercise}

\begin{sole}
    \item Let $p = a_0 + a_1 x + a_2 x^2 + a_3 x^3 + a_4 x^4 \in U$. Then we have
    \[p''(x) = 2a_2 + 6a_3 x + 12a_4 x^2,\]
    and so $p''(6) = 0$ gives us the equation
    \[2a_2 + 36a_3 + 432a_4 = 0.\]
    We can solve this equation for $a_2$ in terms of $a_3, a_4$:
    \[a_2 = -18a_3 - 216a_4.\]
    Then we can write any polynomial in $U$ as a linear combination of the following four polynomials:
    \[u_1 = 1, \quad u_2 = x, \quad u_3 = -18x^2 + x^3, \quad u_4 = -216x^2 + x^4.\]
    Moreover, these four polynomials are linearly independent, since the only way to write the zero polynomial as a linear combination of $u_1, u_2, u_3, u_4$ is to take all coefficients to be zero. Thus, $\set{u_1, u_2, u_3, u_4}$ is a basis of $U$.
    \item We can extend the basis $\set{u_1, u_2, u_3, u_4}$ of $U$ to a basis of $\pp_4(\r)$ by adding one more polynomial that is linearly independent of $u_1, u_2, u_3, u_4$. For example, we can add the polynomial
    \[u_5 = x^2,\]
    which is linearly independent of $u_1, u_2, u_3, u_4$ since it is not in $U$. Thus, $\set{u_1, u_2, u_3, u_4, u_5}$ is a basis of $\pp_4(\r)$.
    \item Let $W = \spn u_5 = \spn x^2$. Then we have
    \[\pp_4(\r) = U + W,\]
    since every polynomial in $\pp_4(\r)$ can be written as a sum of a polynomial in $U$ and a polynomial in $W$. Moreover, if $p \in U \cap W$, then $p = \alpha x^2$ for some $\alpha \in \r$. Since $p''(6) = 0$, we have $2\alpha = 0$, which implies that $\alpha = 0$. Thus, $U \cap W = \set{0}$, and so we have
    \[\pp_4(\r) = U \op W. \qh\]
\end{sole}

\begin{exercise}[2c5]
    \begin{subexercises}
        \item Let $U = \set{p \in \pp_4(\f) \mid p(2) = p(5)}$. Find a basis of $U$.
        \item Extend the basis in (a) to a basis of $\pp_4(\f)$.
        \item Find a subspace $W$ of $\pp_4(\f)$ such that $\pp_4(\f) = U \op W$.
    \end{subexercises}
\end{exercise}

\begin{sole}
    \item Let $p = a_0 + a_1 x + a_2 x^2 + a_3 x^3 + a_4 x^4 \in U$. Then we have $p(2) = p(5)$, which gives us the equation
    \[a_0 + 2a_1 + 4a_2 + 8a_3 + 16a_4 = a_0 + 5a_1 + 25a_2 + 125a_3 + 625a_4.\]
    We can simplify this equation to get
    \[-3a_1 - 21a_2 - 117a_3 - 609a_4 = 0.\]
    We can solve this equation for $a_1$ in terms of $a_2, a_3, a_4$:
    \[a_1 = -7a_2 - 39a_3 - 203a_4.\]
    Then we can write any polynomial in $U$ as a linear combination of the following four polynomials:
    \[u_1 = 1, \quad u_2 = -7x + x^2, \quad u_3 = -39x + x^3, \quad u_4 = -203x + x^4.\]
    Moreover, these four polynomials are linearly independent, since the only way to write the zero polynomial as a linear combination of $u_1, u_2, u_3, u_4$ is to take all coefficients to be zero. Thus, $\set{u_1, u_2, u_3, u_4}$ is a basis of $U$.
    \item We can extend the basis $\set{u_1, u_2, u_3, u_4}$ of $U$ to a basis of $\pp_4(\f)$ by adding one more polynomial that is linearly independent of $u_1, u_2, u_3, u_4$. For example, we can add the polynomial
    \[u_5 = x,\]
    which is linearly independent of $u_1, u_2, u_3, u_4$ since it is not in $U$. Thus, $\set{u_1, u_2, u_3, u_4, u_5}$ is a basis of $\pp_4(\f)$.
    \item Let $W = \spn u_5 = \spn x$. Then we have
    \[\pp_4(\f) = U + W,\]
    since every polynomial in $\pp_4(\f)$ can be written as a sum of a polynomial in $U$ and a polynomial in $W$. Moreover, if $p \in U \cap W$, then $p = \alpha x$ for some $\alpha \in \f$. Since $p(2) = p(5)$, we have $\alpha 2 = \alpha 5$, which implies that $\alpha = 0$. Thus, $U \cap W = \set{0}$, and so we have
    \[\pp_4(\f) = U \op W. \qh\]
\end{sole}

\begin{exercise}[2c6]
    \begin{subexercises}
        \item Let $U = \set{p \in \pp_4(\f) \mid p(2) = p(5) = p(6)}$. Find a basis of $U$.
        \item Extend the basis in (a) to a basis of $\pp_4(\f)$.
        \item Find a subspace $W$ of $\pp_4(\f)$ such that $\pp_4(\f) = U \op W$.
    \end{subexercises}
\end{exercise}

\begin{sole}
    \item Let $p = a_0 + a_1 x + a_2 x^2 + a_3 x^3 + a_4 x^4 \in U$. Then we have $p(2) = p(5) = p(6)$. We know from \autoref{ex:2c5} that the condition $p(2) = p(5)$ gives us the equation
    \[a_1 = -7a_2 - 39a_3 - 203a_4.\]
    The condition $p(5) = p(6)$ gives us the equation
    \[5a_1 + 25a_2 + 125a_3 + 625a_4 = 6a_1 + 36a_2 + 216a_3 + 1296a_4,\]
    which simplifies to
    \[a_1 + 11a_2 + 91a_3 + 671a_4 = 0.\]
    Since we know that $a_1 = -7a_2 - 39a_3 - 203a_4$, we can substitute this into the second equation to get
    \[4a_2 + 52a_3 + 468a_4 = 0.\]
    We can solve this equation for $a_2$ in terms of $a_3, a_4$:
    \[a_2 = -13a_3 - 117a_4.\]
    Then $a_1 = 52a_3 + 616a_4$. Thus, we can write any polynomial in $U$ as a linear combination of the following three polynomials:
    \[u_1 = 1, \quad u_2 = 52x - 13x^2 + x^3, \quad u_3 = 616x - 117x^2 + x^4.\]
    Moreover, these three polynomials are linearly independent, since the only way to write the zero polynomial as a linear combination of $u_1, u_2, u_3$ is to take all coefficients to be zero. Thus, $\set{u_1, u_2, u_3}$ is a basis of $U$.
    \item We can extend the basis $\set{u_1, u_2, u_3}$ of $U$ to a basis of $\pp_4(\f)$ by adding two more polynomials that are linearly independent of $u_1, u_2, u_3$. For example, we can add the polynomials
    \[u_4 = x^2, \quad u_5 = x^3,\]
    which are linearly independent of $u_1, u_2, u_3$ since they are not in $U$. Thus, $\set{u_1, u_2, u_3, u_4, u_5}$ is a basis of $\pp_4(\f)$.
    \item Let $W = \spn(u_4, u_5) = \spn(x^2, x^3)$. Then we have
    \[\pp_4(\f) = U + W,\]
    since every polynomial in $\pp_4(\f)$ can be written as a sum of a polynomial in $U$ and a polynomial in $W$. Moreover, if $p \in U \cap W$, then $p = \alpha x^2 + \beta x^3$ for some $\alpha, \beta \in \f$. Since $p(2) = p(5) = p(6)$, we have
    \[4\alpha + 8\beta = 25\alpha + 125\beta = 36\alpha + 216\beta.\]
    Solving this system of equations, we find that $\alpha = \beta = 0$. Thus, $U \cap W = \set{0}$, and so we have
    \[\pp_4(\f) = U \op W. \qh\]
\end{sole}

\begin{exercise}[2c7]
    \begin{subexercises}
        \item Let
        \[U = \set*{p \in \pp_4(\r) \left|~\int_{-1}^{1} p = 0 \right.}.\]
        Find a basis of $U$.
        \item Extend the basis in (a) to a basis of $\pp_4(\r)$.
        \item Find a subspace $W$ of $\pp_4(\r)$ such that $\pp_4(\r) = U \op W$.
    \end{subexercises}
\end{exercise}

\begin{sole}
    \item Let $p = a_0 + a_1 x + a_2 x^2 + a_3 x^3 + a_4 x^4 \in U$. Then we have
    \begin{align*}
        \int_{-1}^1 p & = \int_{-1}^1 (a_0 + a_1 x + a_2 x^2 + a_3 x^3 + a_4 x^4) \dd x \\
        & = \int_{-1}^1 (a_0 + a_2 x^2 + a_4 x^4) \dd x \\
        & = 2\int_0^1 (a_0 + a_2 x^2 + a_4 x^4) \dd x \\
        & = 2\left(a_0 + \frac{a_2}{3} + \frac{a_4}{5}\right) = 0.
    \end{align*}
    We can solve this equation for $a_4$ in terms of $a_0, a_2$:
    \[a_4 = -5a_2 - 15a_0.\]
    Then we can write any polynomial in $U$ as a linear combination of the following four polynomials:
    \[u_1 = 1 - 15x^4, \quad u_2 = x, \quad u_3 = x^2 - 5x^4, \quad u_4 = x^3.\]
    Moreover, these four polynomials are linearly independent, since the only way to write the zero polynomial as a linear combination of $u_1, u_2, u_3, u_4$ is to take all coefficients to be zero. Thus, $\set{u_1, u_2, u_3, u_4}$ is a basis of $U$.
    \item We can extend the basis $\set{u_1, u_2, u_3, u_4}$ of $U$ to a basis of $\pp_4(\r)$ by adding one more polynomial that is linearly independent of $u_1, u_2, u_3, u_4$. For example, we can add the polynomial
    \[u_5 = x^4,\]
    which is linearly independent of $u_1, u_2, u_3, u_4$ since it is not in $U$. Thus, $\set{u_1, u_2, u_3, u_4, u_5}$ is a basis of $\pp_4(\r)$.
    \item Let $W = \spn u_5 = \spn x^4$. Then we have
    \[\pp_4(\r) = U + W,\]
    since every polynomial in $\pp_4(\r)$ can be written as a sum of a polynomial in $U$ and a polynomial in $W$. Moreover, if $p \in U \cap W$, then $p = \alpha x^4$ for some $\alpha \in \r$. Since the integral of an even function over a symmetric interval is nonzero unless the function is identically zero, we have $\alpha = 0$. Thus, $U \cap W = \set{0}$, and so we have
    \[\pp_4(\r) = U \op W. \qh\]
\end{sole}

\begin{exercise}[2c8]
    Suppose $v_1, \ldots, v_m$ is linearly independent in $V$ and $w \in V$. Prove that
    \[\dim \spn(v_1 + w, \ldots, v_m + w) \geq m - 1.\]
\end{exercise}

\begin{sol}
    Observe that $(v_i + w) - (v_{i + 1} + w) = v_i - v_{i + 1}$ for $1 \leq i \leq m - 1$. Then the list of vectors
    \[v_1 - v_2, v_2 - v_3, \ldots, v_{m-1} - v_m\]
    is contained in $\spn(v_1 + w, \ldots, v_m + w)$. We claim that this list is linearly independent. Suppose that
    \[\alpha_1(v_1 - v_2) + \alpha_2(v_2 - v_3) + \cdots + \alpha_{m - 1}(v_{m - 1} - v_m) = 0\]
    for some $\alpha_1, \ldots, \alpha_{m - 1} \in \f$. Then we can rewrite this equation as
    \[\alpha_1 v_1 + (\alpha_2 - \alpha_1)v_2 + \cdots + (\alpha_{m - 1} - \alpha_{m - 2})v_{m - 1} - \alpha_{m - 1} v_m = 0.\]
    Since $v_1, \ldots, v_m$ is linearly independent, we must have $\alpha_1 = 0$, $\alpha_2 - \alpha_1 = 0$, $\ldots$, $\alpha_{m - 1} - \alpha_{m - 2} = 0$, and $-\alpha_{m - 1} = 0$. This implies that $\alpha_1 = \cdots = \alpha_{m - 1} = 0$. Thus, the list is linearly independent, and so we have
    \[\dim \spn(v_1 + w, \ldots, v_m + w) \geq m - 1. \qh\]
\end{sol}

\begin{exercise}[2c9]
    Suppose $m$ is a positive integer and $p_0, p_1, \ldots, p_m \in \pp(\f)$ are such that each $p_k$ has degree $k$. Prove that $p_0, p_1, \ldots, p_m$ is a basis of $\pp_m(\f)$.
\end{exercise}

\begin{sol}
    Since each $p_k$ has degree $k$, we can write each polynomial as
    \[p_k = a_{k0} + a_{k1}x + \cdots + a_{kk}x^k\]
    where $a_{ij} \in \f$ for all $0 \leq i \leq j \leq k$ and $a_{kk} \neq 0$. We claim that $p_0, p_1, \ldots, p_m$ is linearly independent. Suppose that
    \[\alpha_0 p_0 + \alpha_1 p_1 + \cdots + \alpha_m p_m = 0\]
    for some $\alpha_0, \ldots, \alpha_m \in \f$. Then we can write this equation as
    \[(\alpha_0a_{00} + \alpha_1a_{10} + \cdots + \alpha_ma_{m0}) + (\alpha_1a_{11} + \cdots + \alpha_ma_{m1})x + \cdots + (\alpha_ma_{mm})x^m = 0.\]
    Since the zero polynomial has all coefficients equal to zero, we must have
    \[\alpha_m a_{mm} = 0.\]
    Since $a_{mm} \neq 0$, we must have $\alpha_m = 0$. Next, the coefficient of the $x^{m - 1}$ term gives us
    \[\alpha_{m - 1} a_{m - 1, m -1} + \alpha_m a_{m, m - 1} = 0.\]
    Since $\alpha_m = 0$ and $a_{m - 1, m - 1} \neq 0$, we must have $\alpha_{m - 1} = 0$. Continuing in this way, we find that $\alpha_{m - 2} = \cdots = \alpha_0 = 0$. Thus, the list is linearly independent. Moreover, since $\dim \pp_m(\f) = m + 1$, any linearly independent list of $m + 1$ vectors in $\pp_m(\f)$ is a basis. Therefore, $p_0, p_1, \ldots, p_m$ is a basis of $\pp_m(\f)$.
\end{sol}

\begin{exercise}[2c10]
    Suppose $m$ is a positive integer. For $0 \leq k \leq m$, let
    \[p_k(x) = x^k(1 - x)^{m - k}.\]
    Show that $p_0, \ldots, p_m$ is a basis of $\pp_m(\f)$.
\end{exercise}

\begin{sol}
    We first determine what exponents appear in the polynomial $p_k$. Using the Binomial Theorem, we can expand $(1 - x)^{m - k}$ and obtain the expression for $p_k(x)$ as a polynomial in $x$:
    \[p_k(x) = x^k \sum_{i = 0}^{m - k} \binom{m - k}{i} (-1)^i x^i = \sum_{i = 0}^{m - k} \binom{m - k}{i} (-1)^i x^{k + i}.\]
    Thus, the exponents that appear in $p_k$ are $k, k + 1, \ldots, m$, and we may write $p_k$ as
    \[p_k(x) = a_{kk} x^k + a_{k, k + 1} x^{k + 1} + \cdots + a_{km} x^m\]
    where $a_{kk} = 1$ and $a_{kj} \in \f$ for all $k \leq j \leq m$. We claim that $p_0, p_1, \ldots, p_m$ is linearly independent. Suppose that
    \[\alpha_0 p_0 + \alpha_1 p_1 + \cdots + \alpha_m p_m = 0\]
    for some $\alpha_0, \ldots, \alpha_m \in \f$. Then we can write this equation as
    \[(\alpha_0 a_{00}) + (\alpha_0 a_{01} + \alpha_1 a_{11})x + \cdots + (\alpha_0 a_{0m} + \cdots + \alpha_m a_{mm})x^m = 0.\]
    Since the zero polynomial has all coefficients equal to zero, and $a_{00} \neq 0$ by definition of $p_0$, we must have $\alpha_0 = 0$. Next, the coefficient of the $x$ is $\alpha_0 a_{01} + \alpha_1 a_{11} = 0$. Since $\alpha_0 = 0$ and $a_{11} \neq 0$, we must have $\alpha_1 = 0$. Continuing in this way, we find that $\alpha_2 = \cdots = \alpha_m = 0$. Thus, the list is linearly independent. Moreover, since $\dim \pp_m(\f) = m + 1$, any linearly independent list of $m + 1$ vectors in $\pp_m(\f)$ is a basis. Therefore, $p_0, p_1, \ldots, p_m$ is a basis of $\pp_m(\f)$.
\end{sol}

\begin{exercise}[2c11]
    Suppose $U$ and $W$ are both four-dimensional subspaces of $\c^6$. Prove that there exist two vectors in $U \cap W$ such that neither of these vectors is a scalar multiple of the other.
\end{exercise}

\begin{sol}
    Let $U$ and $W$ be four-dimensional subspaces of $\c^6$. By the dimension formula for the sum of two subspaces, we have
    \[\dim(U + W) = \dim U + \dim W - \dim(U \cap W) = 4 + 4 - \dim(U \cap W) = 8 - \dim(U \cap W).\]
    Since $U + W$ is a subspace of $\c^6$, we have $\dim(U + W) \leq 6$. Thus, we have
    \[8 - \dim(U \cap W) \leq 6,\]
    which implies that $\dim(U \cap W) \geq 2$. Therefore, there exist at least two linearly independent vectors in $U \cap W$, and neither of these vectors is a scalar multiple of the other.
\end{sol}

\begin{exercise}[2c12]
    Suppose that $U$ and $W$ are subspaces of $\r^8$ such that $\dim U = 3$, $\dim W = 5$, and $U + W = \r^8$. Prove that $\r^8 = U \op W$.
\end{exercise}

\begin{sol}
    By the dimension formula for the sum of two subspaces, we have
    \[\dim(U + W) = \dim U + \dim W - \dim(U \cap W) = 3 + 5 - \dim(U \cap W) = 8 - \dim(U \cap W).\]
    Since $U + W = \r^8$, we have $\dim(U + W) = 8$. Thus, we have
    \[8 - \dim(U \cap W) = 8,\]
    which implies that $\dim(U \cap W) = 0$. Therefore, $U \cap W = \set{0}$, and so we have
    \[\r^8 = U \op W. \qh\]
\end{sol}

\begin{exercise}[2c13]
    Suppose $U$ and $W$ are both five-dimensional subspaces of $\r^8$. Prove that $U \cap W \neq \set{0}$.
\end{exercise}

\begin{sol}
    By the dimension formula for the sum of two subspaces, we have
    \[\dim(U + W) = \dim U + \dim W - \dim(U \cap W) = 5 + 5 - \dim(U \cap W) = 10 - \dim(U \cap W).\]
    Since $U + W$ is a subspace of $\r^8$, we have $\dim(U + W) \leq 8$. Thus, we have
    \[10 - \dim(U \cap W) \leq 8,\]
    which implies that $\dim(U \cap W) \geq 2$. Therefore, $U \cap W$ is nontrivial, and so we have $U \cap W \neq \set{0}$.
\end{sol}

\begin{exercise}[2c14]
    Suppose $V$ is a ten-dimensional vector space and $V_1, V_2, V_3$ are subspaces of $V$ with $\dim V_1 = \dim V_2 = \dim V_3 = 7$. Prove that $V_1 \cap V_2 \cap V_3 \neq \set{0}$.
\end{exercise}

\begin{sol}
    Let $W = V_1 \cap V_2$. By the dimension formula for the sum of two subspaces, we have
    \[\dim(V_1 + V_2) = \dim V_1 + \dim V_2 - \dim W = 7 + 7 - \dim W = 14 - \dim W.\]
    Since $V_1 + V_2$ is a subspace of $V$, we have $\dim(V_1 + V_2) \leq 10$. Thus, we have $14 - \dim W \leq 10$, which implies that $\dim W \geq 4$. Therefore, $V_1 \cap V_2$ is at least four-dimensional. Now, consider the subspace $W \cap V_3 = V_1 \cap V_2 \cap V_3$. By the dimension formula for the sum of two subspaces, we have
    \[\dim(W + V_3) = \dim W + \dim V_3 - \dim(W \cap V_3) = \dim W + 7 - \dim(W \cap V_3).\]
    Since $W + V_3$ is a subspace of $V$, we have $\dim(W + V_3) \leq 10$. Thus, we have
    \[\dim W + 7 - \dim(W \cap V_3) \leq 10,\]
    which implies that $\dim(W \cap V_3) \geq \dim W + 7 - 10 = \dim W - 3$. Since $\dim W \geq 4$, we have $\dim(W \cap V_3) \geq 1$. Therefore, $V_1 \cap V_2 \cap V_3$ is nontrivial, and so we have $V_1 \cap V_2 \cap V_3 \neq \set{0}$.
\end{sol}

\begin{exercise}[2c15]
    Suppose $V$ is finite-dimensional and $V_1, V_2, V_3$ are subspaces of $V$ with $\dim V_1 + \dim V_2 + \dim V_3 > 2 \dim V$. Prove that $V_1 \cap V_2 \cap V_3 \neq \set{0}$.
\end{exercise}

\begin{sol}
    Let $W = V_1 \cap V_2$. By the dimension formula for the sum of two subspaces, we have
    \[\dim(V_1 + V_2) = \dim V_1 + \dim V_2 - \dim W.\]
    Since $V_1 + V_2$ is a subspace of $V$, we have $\dim(V_1 + V_2) \leq \dim V$. Thus, we have
    \[\dim V_1 + \dim V_2 - \dim W \leq \dim V,\]
    which implies that $\dim W \geq \dim V_1 + \dim V_2 - \dim V$. Now, consider the subspace $W \cap V_3 = V_1 \cap V_2 \cap V_3$. By the dimension formula for the sum of two subspaces, we have
    \[\dim(W + V_3) = \dim W + \dim V_3 - \dim(W \cap V_3).\]
    Since $W + V_3$ is a subspace of $V$, we have $\dim(W + V_3) \leq \dim V$. Thus, we have
    \[\dim W + \dim V_3 - \dim(W \cap V_3) \leq \dim V,\]
    which implies that $\dim(W \cap V_3) \geq \dim W + \dim V_3 - \dim V$. Since $\dim W \geq \dim V_1 + \dim V_2 - \dim V$, we have
    \[\dim(W \cap V_3) \geq (\dim V_1 + \dim V_2 - \dim V) + \dim V_3 - \dim V = \dim V_1 + \dim V_2 + \dim V_3 - 2 \dim V.\]
    Since $\dim V_1 + \dim V_2 + \dim V_3 > 2 \dim V$, we have $\dim(W \cap V_3) > 0$. Therefore, $V_1 \cap V_2 \cap V_3$ is nontrivial, and so we have $V_1 \cap V_2 \cap V_3 \neq \set{0}$.
\end{sol}

\begin{exercise}[2c16]
    Suppose $V$ is finite-dimensional and $U$ is a subspace of $V$ with $U \neq V$. Let $n = \dim V$ and $m = \dim U$. Prove that there exist $n - m$ subspaces of $V$, each of dimension $n - 1$, whose intersection equals $U$.
\end{exercise}

\begin{sol}
    Let $U$ be a subspace of $V$ with $\dim U = m < n = \dim V$. We can extend a basis of $U$ to a basis of $V$. Let $\set{u_1, \ldots, u_m}$ be a basis of $U$, and extend it to a basis of $V$ by adding vectors $v_1, \ldots, v_{n - m}$. Then $\set{u_1, \ldots, u_m, v_1, \ldots, v_{n - m}}$ is a basis of $V$. For each $i = 1, \ldots, n - m$, let
    \[W_i = \spn(u_1, \ldots, u_m, v_1, \ldots, v_{i - 1}, v_{i + 1}, \ldots, v_{n - m}).\]
    Then each $W_i$ is a subspace of $V$ with dimension $n - 1$, since it is spanned by $n - 1$ linearly independent vectors. Moreover, we have
    \[U = W_1 \cap W_2 \cap \cdots \cap W_{n - m},\]
    since any vector in the intersection must be a linear combination of the basis vectors of $U$, and any vector in $U$ is contained in all the subspaces $W_i$. Therefore, there exist $n - m$ subspaces of $V$, each of dimension $n - 1$, whose intersection equals $U$.
\end{sol}

\begin{exercise}[2c17]
    Suppose that $V_1, \ldots, V_m$ are finite-dimensional subspaces of $V$. Prove that $V_1 + \cdots + V_m$ is finite-dimensional and
    \[\dim(V_1 + \cdots + V_m) \leq \dim V_1 + \cdots + \dim V_m.\]
\end{exercise}

\begin{sol}
    We will prove the statement by induction on $m$. For the base case, let $m = 2$. Then we have
    \[V_1 + V_2 = V_1 + V_2,\]
    which is finite-dimensional and
    \[\dim(V_1 + V_2) = \dim V_1 + \dim V_2 - \dim(V_1 \cap V_2) \leq \dim V_1 + \dim V_2.\]
    Hence, the statement holds for $m = 2$. Now, suppose that the statement holds for some positive integer $m$. That is, suppose that if $V_1, \ldots, V_m$ are finite-dimensional subspaces of $V$, then $V_1 + \cdots + V_m$ is finite-dimensional and
    \[\dim(V_1 + \cdots + V_m) \leq \dim V_1 + \cdots + \dim V_m.\]
    We want to show that the statement holds for $m + 1$. Let $V_1, \ldots, V_{m + 1}$ be finite-dimensional subspaces of $V$. By the induction hypothesis, we know that $W = V_1 + \cdots + V_m$ is finite-dimensional and
    \[\dim W \leq \dim V_1 + \cdots + \dim V_m.\]
    Then we have
    \[\dim(W + V_{m + 1}) = \dim W + \dim V_{m + 1} - \dim(W \cap V_{m + 1}).\]
    Since $W \cap V_{m + 1}$ is a subspace of $V$, we have $\dim(W \cap V_{m + 1}) \geq 0$. Thus, we have
    \[\dim(W + V_{m + 1}) \leq \dim W + \dim V_{m + 1} \leq \dim V_1 + \cdots + \dim V_m + \dim V_{m + 1}.\]
    Therefore, $V_1 + \cdots + V_{m + 1}$ is finite-dimensional and
    \[\dim(V_1 + \cdots + V_{m + 1}) \leq \dim V_1 + \cdots + \dim V_m + \dim V_{m + 1}. \qh\]
\end{sol}

\begin{exercise}[2c18]
    Suppose $V$ is finite-dimensional, with $\dim V = n \geq 1$. Prove that there exist one-dimensional subspaces $V_1, \ldots, V_n$ of $V$ such that
    \[V = V_1 \op \cdots \op V_n.\]
\end{exercise}

\begin{sol}
    Let $v_1, \ldots, v_n$ be a basis of $V$. For each $i = 1, \ldots, n$, let $V_i = \spn v_i$. Then each $V_i$ is a one-dimensional subspace of $V$, since it is spanned by a single nonzero vector. Moreover, we have
    \[V = V_1 + \cdots + V_n,\]
    and the sum is direct because the basis vectors are linearly independent. Therefore, we have
    \[V = V_1 \op \cdots \op V_n. \qh\]
\end{sol}

\begin{exercise}[2c19]
    Explain why you might guess, motivated by analogy with the formula for the number of elements in the union of three finite sets, that if $V_1, V_2, V_3$ are subspaces of a finite-dimensional vector space, then
    \begin{align*}
        \dim(V_1 + V_2 + V_3) = & \dim V_1 + \dim V_2 + \dim V_3 \\
        & - \dim(V_1 \cap V_2) - \dim(V_1 \cap V_3) - \dim(V_2 \cap V_3) \\
        & + \dim(V_1 \cap V_2 \cap V_3).
    \end{align*}
    Then either prove the formula above or give a counterexample.
\end{exercise}

\begin{sol}
    One may guess that the given formula above is true by analogy with the formula for the number of elements in the union of three finite sets. However, this erroneously conflates that the idea of subspace sums is equivalent to the idea of set unions. In fact, the formula is not true in general. For example, let $V = \r^2$, and let
    \[V_1 = \spn((1, 0)), \quad V_2 = \spn((0, 1)), \quad V_3 = \spn((1, 1)).\]
    It is easy to see that $\dim V_1 = \dim V_2 = \dim V_3 = 1$, but $\dim(V_1 \cap V_2) = \dim(V_1 \cap V_3) = \dim(V_2 \cap V_3) = 0$ and $\dim(V_1 \cap V_2 \cap V_3) = 0$. Thus, the right-hand side of the formula is
    \[1 + 1 + 1 - 0 - 0 - 0 + 0 = 3,\]
    but the left-hand side is $\dim(V_1 + V_2 + V_3) = \dim \r^2 = 2$. Therefore, the formula is not true in general.
\end{sol}

\begin{exercise}[2c20]
    Prove that if $V_1, V_2$, and $V_3$ are subspaces of a finite-dimensional vector space, then
    \begin{align*}
        \dim(V_1 + V_2 + V_3) = & \dim V_1 + \dim V_2 + \dim V_3 \\
        &- \frac{\dim(V_1 \cap V_2) + \dim(V_1 \cap V_3) + \dim(V_2 \cap V_3)}{3} \\
        &- \frac{\dim((V_1 + V_2) \cap V_3) + \dim((V_1 + V_3) \cap V_2) + \dim((V_2 + V_3) \cap V_1)}{3}.
    \end{align*}
\end{exercise}

\begin{sol}
    Note that we need to account for all pairs of subspaces all pairs of sums of subspaces, since the sum of subspaces is not equivalent to the union of sets. We can use the dimension formula repeatedly to obtain the answer. To that end, define the following subspaces:
    \[W_1 = V_1 + V_2, \quad W_2 = V_1 + V_3, \quad W_3 = V_2 + V_3.\]
    Applying the dimension formula to each of $W_1, W_2, W_3$, we have
    \begin{align*}
        \dim W_1 & = \dim V_1 + \dim V_2 - \dim(V_1 \cap V_2), \\
        \dim W_2 & = \dim V_1 + \dim V_3 - \dim(V_1 \cap V_3), \\
        \dim W_3 & = \dim V_2 + \dim V_3 - \dim(V_2 \cap V_3).
    \end{align*}
    Next, we can apply the dimension formula to each of the sums $W_1 + V_3, W_2 + V_2, W_3 + V_1$:
    \begin{align*}
        \dim(W_1 + V_3) & = \dim W_1 + \dim V_3 - \dim(W_1 \cap V_3), \\
        & = \dim V_1 + \dim V_2 - \dim(V_1 \cap V_2) + \dim V_3 - \dim(W_1 \cap V_3) \\
        \dim(W_2 + V_2) & = \dim W_2 + \dim V_2 - \dim(W_2 \cap V_2), \\
        & = \dim V_1 + \dim V_3 - \dim(V_1 \cap V_3) + \dim V_2 - \dim(W_2 \cap V_2) \\
        \dim(W_3 + V_1) & = \dim W_3 + \dim V_1 - \dim(W_3 \cap V_1), \\
        & = \dim V_2 + \dim V_3 - \dim(V_2 \cap V_3) + \dim V_1 - \dim(W_3 \cap V_1).
    \end{align*}
    Note that $W_1 + V_3 = W_2 + V_2 = W_3 + V_1 = V_1 + V_2 + V_3$. Thus, we can add the three equations above and divide by $3$ to obtain the desired formula:
    \begin{align*}
        \dim(V_1 + V_2 + V_3) = & \dim V_1 + \dim V_2 + \dim V_3 \\
        &- \frac{\dim(V_1 \cap V_2) + \dim(V_1 \cap V_3) + \dim(V_2 \cap V_3)}{3} \\
        &- \frac{\dim((V_1 + V_2) \cap V_3) + \dim((V_1 + V_3) \cap V_2) + \dim((V_2 + V_3) \cap V_1)}{3}. \qh
    \end{align*}
\end{sol}